\section{Deferred Proofs}
\label{appx:proof}
\label{appx:auxiliary-proofs}

This appendix proves the claims in Sections~\ref{sec:blocks}
and~\ref{sec:together} and the packed PCS claim in
Appendix~\ref{appx:pcs}, using the notation of each statement.

\subsection{Proof of Proposition~\ref{proposition:ip-complexity}}
\label{appx:proof-ip}

\begin{proof}
\noindent\textbf{Completeness.}
Expand the Laurent polynomial in Eq.~\eqref{eq:ip-laurent-polynomial}:
\[
    f(X)=
    \sum_{\tau=-(\kappa-1)}^{\kappa-1}f_\tau X^\tau,
    \qquad
    f_\tau=
    \sum_{\substack{j,j'\in[\kappa]\\j-j'=\tau}}
       \langle\bm a^{(j)},\bm b^{(j')}\rangle.
\]
Its constant coefficient is \(f_0=\langle\bm a,\bm b\rangle\).
Thus the auxiliary polynomials in Algorithm~\ref{alg:ip}
satisfy Eq.~\eqref{eq:ip-correction-identity} for a correct
inner-product claim. Evaluation at \(\ext x\neq0\) gives the
length-\(m\) claim in
Eq.~\eqref{eq:packed-inner-product-check}.

The terminal values of Fact~\ref{fact:ip} have the matrix form:
\[
\begin{aligned}
    \ext v_1
      &=\tensorspan{\ext{\bm y}}^\top\bm A\ext{\bm\rho}_A,
    &\ext{\bm\rho}_A
      &=(1,\ext x,\ldots,\ext x^{\kappa-1})^\top,\\
    \ext v_2
      &=\tensorspan{\ext{\bm y}}^\top\bm B\ext{\bm\rho}_B,
    &\ext{\bm\rho}_B
      &=(1,\ext x^{-1},\ldots,\ext x^{-(\kappa-1)})^\top.
\end{aligned}
\]
Both have tensor-structured row weights and fit the
packed-opening relation in Eq.~\eqref{eq:pcs-opening-relation}.

\smallskip\noindent\textbf{Soundness.}
Suppose \(c\neq\langle\bm a,\bm b\rangle\), and fix the
degree-\((\kappa-2)\) polynomials \(h_+,h_-\) before sampling
\(\ext x\). The Laurent residual
\(f(X)-c-Xh_+(X)-X^{-1}h_-(X^{-1})\) has a nonzero
constant coefficient. Multiplying by
\(X^{\kappa-1}\) gives a nonzero polynomial of degree at most
\(2\kappa-2\), so the evaluation check accepts with probability
at most \((2\kappa-2)/(|\ExtensionField|-1)\).
Conditioned on a false evaluated claim, Fact~\ref{fact:ip}
contributes at most \(2\log_2m/|\ExtensionField|\).
A union bound gives the stated error, excluding the terminal
packed openings.

\smallskip\noindent\textbf{Efficiency.}
Each ordered pair of distinct columns contributes a
length-\(m\) base-field inner product to the correction
coefficients, for \(m\kappa(\kappa-1)=N(\kappa-1)\) base-field
multiplications. Given the powers of \(\ext x\) and
\(\ext x^{-1}\), packing both matrices uses \(2N\,\be\).
Generating these powers and evaluating the corrections costs
\(O(\kappa)\,\ee\); Fact~\ref{fact:ip} adds
\((4m+O(\log m))\,\ee\).
Thus the total arithmetic cost is:
\[
    O(N\kappa)\,\bb+O(N)\,\be+O(m+\kappa)\,\ee
       =O(N\kappa+\kappa^3)\,\bb.
\]
Streaming the cross-column inner products retains only
\(O(\kappa)\) base-field correction coefficients.
The witness matrices and packed vectors use \(O(N)\)
base-field coordinates, and the challenge powers use
\(O(\kappa^2)\). An in-place inner-product reduction therefore
has peak memory \(O(N+\kappa^2)\Field\).
\end{proof}

\subsection{Proof of Proposition~\ref{proposition:hadamard-complexity}}
\label{appx:proof-hadamard}

\begin{proof}
\noindent\textbf{Completeness.}
The honest column message is
\(\ext{\bm c}^{\top}
=\ext{\bm l}^{\top}(\bm A\circ\bm B)\), so the scalar check
\(\ext c=\langle\ext{\bm c},\ext{\bm r}\rangle\) accepts.
At every interpolation node, the residual satisfies:
\[
    f(\gamma_j)
      =\langle\ext{\bm l},
          \bm a^{(j)}\circ\bm b^{(j)}\rangle-\ext c_j
      =0,
    \qquad j\in[\kappa].
\]
Hence \(f=\Psi h\) for a polynomial \(h\) of degree at most
\(\kappa-2\). Evaluating this identity at \(\ext x\) gives:
\[
    \langle\ext{\bm a},\ext{\bm b}\circ\ext{\bm l}\rangle
      =\ext c_x+\Psi(\ext x)h(\ext x).
\]
The terminal values of Fact~\ref{fact:ip} satisfy:
\[
\begin{aligned}
    \ext v_1
      &=\tensorspan{\ext{\bm y}}^\top
          \bm A\ext{\bm\rho}(\ext x),\\
    \ext v_2
      &=\bigl(\ext{\bm l}\circ\tensorspan{\ext{\bm y}}\bigr)^\top
          \bm B\ext{\bm\rho}(\ext x).
\end{aligned}
\]
Writing the supplied row weight as
\(\ext{\bm l}=\bigotimes_{j=1}^{\log_2m}
(\ext l_{j,0},\ext l_{j,1})^\top\), its product with the MLE
weight factors as:
\[
    \ext{\bm l}\circ\tensorspan{\ext{\bm y}}
      =\bigotimes_{j=1}^{\log_2m}
        \begin{pmatrix}
            \ext l_{j,0}(1-\ext y_j)\\
            \ext l_{j,1}\ext y_j
        \end{pmatrix}.
\]
Both terminal claims fit Eq.~\eqref{eq:pcs-opening-relation},
with Lagrange column weights and tensor-structured row weights.

\smallskip\noindent\textbf{Soundness.}
Suppose the weighted Hadamard claim is false but the scalar
check accepts. Then the column message differs from
\(\ext{\bm l}^{\top}(\bm A\circ\bm B)\) in at least one
coordinate, and \(f(\gamma_j)\neq0\) for some \(j\).
For every degree-\((\kappa-2)\) polynomial \(h\) fixed before
\(\ext x\), the residual \(f-\Psi h\) is therefore nonzero
and has degree at most \(2\kappa-2\).
A random \(\ext x\in\ExtensionField\setminus\mathcal H\)
makes this residual vanish with probability at most
\((2\kappa-2)/(|\ExtensionField|-\kappa)\).
Fact~\ref{fact:ip} contributes at most
\(2\log_2m/|\ExtensionField|\), giving:
\[
    \epsilon_{\mathsf{Had}}
      \leq
      \frac{2\kappa-2}{|\ExtensionField|-\kappa}
      +\frac{2\log_2m}{|\ExtensionField|}
      \leq
      \frac{2\kappa-2}{|\ExtensionField|-\kappa}
      +\frac{2\log_2N}{|\ExtensionField|}.
\]
The first inequality bounds Algorithm~\ref{alg:hadamard};
the second gives the stated bound. Terminal packed openings
are accounted for separately.
The argument applies unchanged when the interpolation nodes
lie in \(\mathbb K\subseteq\ExtensionField\).

\smallskip\noindent\textbf{Efficiency over the base field.}
First suppose \(\mathbb K=\Field\).
Computing \(\ext{\bm c}\) uses \(O(N)\,\bb+O(N)\,\be\).
At each of the \(\kappa-1\) additional interpolation nodes,
evaluating both witness encodings takes \(O(N)\,\bb\).
Pointwise multiplication and contraction with \(\ext{\bm l}\)
then take \(O(m)\,\bb+O(m)\,\be\).
Over all nodes, these operations cost
\(O(N\kappa)\,\bb+O(N)\,\be\).
Evaluating \(g_C\), interpolating the residual, and dividing by
\(\Psi\) contribute \(O(\kappa^2)\,\be\).

After \(\ext x\) is sampled, packing both matrices costs
\(O(N)\,\be\).
Expanding the row factors, forming
\(\ext{\bm b}\circ\ext{\bm l}\), and running Fact~\ref{fact:ip}
cost \(O(m)\,\ee\).
The terminal row factors require only componentwise
multiplication of their two-entry representations.
With precomputed Lagrange coefficients, the public query
and polynomial evaluations cost
\(O(\kappa^2)\,\be+O(\kappa)\,\ee\).
The total is therefore:
\[
    O(N\kappa)\,\bb
      +O(N+\kappa^2)\,\be
      +O(m+\kappa)\,\ee
      =O(N\kappa+\kappa^3)\,\bb.
\]

\smallskip\noindent\textbf{Intermediate-field tables.}
Let \([\mathbb K:\Field]=t>1\), with \(t\mid\kappa\) and
\(|\mathbb K|\geq2\kappa-1\).
Use the addition and multiplication tables specified in
Section~\ref{sec:hadamard}, charging each lookup as \(O(1)\).
Evaluating the witness encodings at all additional nodes
requires \(O(N\kappa)\) table operations.
In \(\mathbb K\)-coordinates, each extension-field element
has length \(\kappa/t\), so weighting and accumulating the
\(O(N)\) local products uses \(O(N\kappa/t)\) lookups.
The \(O(\kappa^2)\) mixed operations for \(g_C\), interpolation,
and division cost \(O(\kappa^3/t)\).
The remaining packing and inner-product steps retain their
\(O(N\kappa+\kappa^3)\,\bb\) bound. Thus the same online
time bound holds, excluding construction of the fixed tables.

\smallskip\noindent\textbf{Memory.}
Processing rows and sample points sequentially retains
\(O(\kappa)\) extension-field residual values, followed by
\(O(\kappa)\) extension-field coefficients after interpolation.
The packed vectors and inner-product reduction use \(O(m)\)
extension-field elements. Including the witness matrices,
this is \(O(N+\kappa^2)\) base-field coordinates.
The fixed interpolation data use \(O(\kappa^2t)\)
coordinates, and the tables use \(O(tq^{2t})\) when \(t>1\);
no field tables are needed when \(t=1\).
The peak memory is therefore
\(O\bigl(N+\kappa^2t+(t-1)q^{2t}\bigr)\Field\).
\end{proof}

\subsection{Proof of Proposition~\ref{proposition:permutation-complexity}}
\label{appx:proof-permutation}

\begin{proof}
\noindent\textbf{Completeness.}
The coefficient tree in
Eq.~\eqref{eq:packed-permutation-product-tree} computes the
products of \(2^\tau\) shifted witness entries in each column.
Evaluating its last layer at \(\ext z\) gives \(m\)
extension-field entries, whose product is
\(\mathsf{GrandProd}(\ext z\bm1+\bm a)\).
The two grand products agree whenever \(\bm a\) and \(\bm b\)
contain the same multiset of entries.

For a transition from level \(\tau+1\) to level \(\tau\),
Lagrange interpolation gives:
\[
    \bm g_{\tau,U}(Z,\gamma_j)
      \circ\bm g_{\tau,D}(Z,\gamma_j)
       =\bm C_{\tau+1}(Z)_{*,j},
    \qquad j\in[\kappa].
\]
For the honest column claims
\(\ext{\bm u}_{\tau+1}^\top
=\ext{\bm l}_{\tau+1}^\top\bm C_{\tau+1}(\ext z)\),
the residual \(f(X)\) therefore vanishes on \(\mathcal H\).
Since its degree is at most \(2\kappa-2\), it satisfies
\(f(X)=\Psi(X)h(X)\) with \(\deg h\leq\kappa-2\).
At an additional node, coefficient aggregation commutes with
evaluation at \(\ext z\):
\[
\begin{aligned}
    F_{\tau+1,j}(\ext z)
      &=\sum_{\nu=0}^{2^{\tau+1}}
        \langle\ext{\bm l}_{\tau+1},
                 \bm p_{\tau+1,j,\nu}\rangle\ext z^\nu\\
      &=\left\langle\ext{\bm l}_{\tau+1},
             \bm g_{\tau,U}(\ext z,\gamma'_j)
               \circ\bm g_{\tau,D}(\ext z,\gamma'_j)
        \right\rangle.
\end{aligned}
\]
Algorithm~\ref{alg:permutation} therefore recovers \(h\) from
\(\kappa-1\) samples. Reencoding \(\bm C_\tau(Z)\) at each
level keeps the \(X\)-degree of \(\bm g_\tau(Z,X)\) at most
\(\kappa-1\).

Each evaluated relation is proved by weighted sumcheck.
For length-\(n\) child vectors and a tensor weight
\(\ext{\bm l}\), the sumcheck polynomial is:
\begin{equation}
    \widetilde{\ext{\bm l}}(\bm Y)\,
    \widetilde{\ext{\bm u}}(\bm Y)\,
    \widetilde{\ext{\bm v}}(\bm Y).
\label{eq:permutation-hadamard-as-ip}
\end{equation}
It has individual degree at most three.
The two child claims are evaluated at the same point
\(\ext{\bm y}\). Their affine combination with a fresh
\(\ext y_0\) is the MLE claim on the concatenated parent
vector at \(\ext y_0\mathbin\|\ext{\bm y}\).
This preserves the recursive matrix claim
\(\ext v_\tau=\ext{\bm l}_\tau^\top
\bm C_\tau(\ext z)\ext{\bm r}_\tau\).
At level zero, the row MLE weights and the column Lagrange
weights both sum to one. Removing the shift leaves:
\[
    \ext v_0-\ext z
       =\ext{\bm l}_0^\top\bm A\ext{\bm r}_0.
\]
This is the packed-opening interface of
Eq.~\eqref{eq:pcs-opening-relation}. The same argument applies
to \(\bm B\).

\smallskip\noindent\textbf{Soundness.}
For a weighted sumcheck on length-\(n\) vectors,
Eq.~\eqref{eq:permutation-hadamard-as-ip} contributes at most
\(3\log_2n/|\ExtensionField|\).
If at least one of the two terminal child claims is false,
their errors form a nonzero affine polynomial in the fresh
merge challenge \(\ext y_0\). The merged claim is then correct
with probability at most \(1/|\ExtensionField|\).
A transition thus preserves a false claim except with
probability at most \((3\log_2n+1)/|\ExtensionField|\).
The extension-field tree has child lengths
\(1,2,\ldots,m/2\). Summing its errors gives:
\[
    \sum_{i=0}^{\log_2m-1}
       \frac{3i+1}{|\ExtensionField|}
       =\frac{\frac32\log_2m(\log_2m-1)+\log_2m}
                    {|\ExtensionField|}.
\]

At coefficient level \(\tau+1\), a false incoming matrix
claim either fails the scalar check against
\(\ext{\bm u}_{\tau+1}\), or makes at least one column claim
false. In the latter case \(f(X)\) does not vanish on all of
\(\mathcal H\), and \(f(X)-\Psi(X)h(X)\) is nonzero for every
claimed \(h\) of degree at most \(\kappa-2\).
Since the message is fixed before \(\ext x\) is sampled,
this polynomial vanishes at \(\ext x\) with probability at
most \((2\kappa-2)/(|\ExtensionField|-\kappa)\).
The subsequent weighted sumcheck and merge have child length
\(m/2^{\tau+1}\) and error at most
\((3\log_2(m/2^{\tau+1})+1)/|\ExtensionField|\).
Summing over \(\tau=0,\ldots,\log_2\kappa-1\) and adding the
extension-field-tree error gives the stated
\(\epsilon_{\mathsf{GP}}\). This bound is conditional on
authentication of the terminal packed claim.

If the committed vectors are not permutations of one another,
their fingerprint difference is the nonzero polynomial:
\[
    \prod_{i=1}^N(Z+a_i)-\prod_{i=1}^N(Z+b_i).
\]
Its degree is at most \(N-1\), so the true fingerprints
coincide with probability at most
\(\min\{1,(N-1)/|\ExtensionField|\}\).
Otherwise, at least one claimed grand product is false.
A union bound over the two reductions and their terminal
authentication gives:
\[
    \epsilon_{\mathsf{Perm}}
       \leq\min\left\{1,\frac{N-1}{|\ExtensionField|}\right\}
          +2\epsilon_{\mathsf{GP}}+\epsilon_{\mathsf{PCS}}.
\]

\smallskip\noindent\textbf{Efficiency over \(\Field\).}
At level \(\tau\in[\log_2\kappa]\), the coefficient matrix
of \(\bm C_{\tau-1}(Z)\) has \(O(m)\) row-degree pairs.
Evaluating its \(\kappa\)-column encodings at \(\kappa-1\)
additional nodes costs \(O(N\kappa)\,\bb\).
Multiplying the resulting degree-\(2^{\tau-1}\) polynomials
costs \(O(m\kappa2^\tau)\,\bb\).
These products, including construction of the coefficient
tree, sum to \(O(N\kappa)\,\bb\).

The off-domain products and column claims contain \(O(N)\)
coefficients per level, whose aggregation costs \(O(N)\,\be\).
There are \(O(\kappa)\) aggregated polynomials of degree
at most \(2^\tau\), so their subsequent evaluation costs:
\[
    \sum_{\tau=1}^{\log_2\kappa}
        O(\kappa2^\tau)\,\ee
       =O(\kappa^2)\,\ee.
\]
The shared powers of \(\ext z\) require
\(O(\kappa)\,\ee\) once.
Forming the vectors for the short weighted sumchecks costs
\(O(N)\,\be+O(m)\,\ee\) per coefficient level.
Their total sumcheck length, together with the
extension-field tree, is \(O(m)\).
Interpolation and public queries contribute
\(O(\kappa^3)\,\bb\) per level.
Using \(m\geq\kappa\), the total is:
\[
    O(N\kappa\log_2\kappa+\kappa^4)\,\bb.
\]
The \(\kappa^4\) term accounts for evaluation of the
aggregated extension-field polynomials.

\smallskip\noindent\textbf{Intermediate-field implementation.}
For \(q<2\kappa-1\), use the addition and multiplication tables
of Section~\ref{sec:hadamard} over \(\mathbb K\), with the same
unit-cost lookup convention.
Precomputed Lagrange evaluations allow the
\(O(N\kappa)\) encoding operations at each level to be
performed over \(\mathbb K\).
The row-polynomial products likewise use
\(O(N\kappa)\) table operations in total.
The coefficients of \(\bm C_\tau(Z)\) remain in \(\Field\);
only the off-domain encodings and products use \(\mathbb K\).

In a compatible tower basis, a mixed
\(\mathbb K\times\ExtensionField\) product and its
accumulation cost \(O(\kappa)\) in base-field coordinates.
Schoolbook multiplication over \(\mathbb K\), using the
tables for coordinate products, costs
\(O(\kappa^2/t)\) per extension-field multiplication.
The preceding count becomes
\(O\bigl(N\kappa\log_2\kappa+\kappa^4/t\bigr)\,\bb\),
which implies the coarser bound in the proposition.

\smallskip\noindent\textbf{Space.}
One coefficient-tree level contains \(O(N)\) base-field
coefficients. The levels can be reconstructed from
\(\bm C_0(Z)\) during the reverse traversal:
reconstructing level \(\tau\) costs \(O(N2^\tau)\,\bb\),
and these costs sum to \(O(N\kappa)\,\bb\).
Thus the prover need not retain all coefficient levels.
Processing one additional node at a time uses
\(O(mt)\subseteq O(N)\) coordinates for its row polynomials.
The aggregated polynomials use \(O(\kappa^2)\)
extension-field coefficients, or \(O(\kappa^3)\)
base-field coordinates. The extension-field tree and
sumcheck tables use \(O(m\kappa)=O(N)\) coordinates.
Including interpolation data and the intermediate-field
tables, peak memory is \(O(N+\kappa^3+tq^{2t})\)
base-field coordinates.
When \(q\geq2\kappa-1\), no field tables are needed, giving
\(O(N+\kappa^3)\) base-field coordinates.
The time bounds exclude preprocessing, commitments, and
terminal openings.
\end{proof}

\subsection{Proof of Proposition~\ref{proposition:packed-pcs}}
\label{appx:proof-pcs}

\begin{proof}
\noindent\textbf{Binding and completeness.}
Each entry of \(\bm A\ext{\bm\beta}\) uniquely determines
the corresponding row of \(\bm A\). This injective encoding
preserves binding and, when provided by the underlying PCS,
commitment hiding.

For an honest prover, \(\bm U=\bm\Xi^\top\bm A\), so
Eq.~\eqref{eq:pcs-encoded-matrix} and its random projection hold.
We show that the verifier also computes the public terminal value
exactly.

By Eq.~\eqref{eq:pcs-multiplication-matrix}, multiplication by
\(\ext z\) acts on a column of base-field coordinates through
\(\bm M_{\ext z}\). This map is \(\Field\)-linear in \(\ext z\),
and \(\bm M_{1}=\bm I_\kappa\).
We prove by backward induction that
\(\bm W_j=\bm Y_j^\top\bm\Xi_j\).
Since \(\ext\beta_1=1\), the empty suffix has coordinate row
\(\bm e_1^\top\), giving
\(\bm Y_{\ell+1}^\top\bm\Xi_{\ell+1}
=\bm e_1\bm e_1^\top=\bm W_{\ell+1}\).
For the induction step, apply the block identities in
Eq.~\eqref{eq:pcs-query-blocks} and the induction hypothesis:
\[
\begin{aligned}
    \bm Y_j^\top\bm\Xi_j
    &=\bm M_{1-\ext y_j}\bm W_{j+1}
        \bm M_{\ext\xi^{(j)}_0}^\top\\
    &\quad{}+\bm M_{\ext y_j}\bm W_{j+1}
        \bm M_{\ext\xi^{(j)}_1}^\top\\
    &=\bm W_{j+1}\bm M_{\ext\xi^{(j)}_0}^\top\\
    &\quad{}+\bm M_{\ext y_j}\bm W_{j+1}
        \bm M_{\ext\xi^{(j)}_1-\ext\xi^{(j)}_0}^\top\\
    &=\bm W_j.
\end{aligned}
\]
The second equality uses linearity of \(\bm M\); the last
is Eq.~\eqref{eq:pcs-query-recurrence}.
Taking \(j=1\) gives \(\bm W=\bm W_1=\bm Y^\top\bm\Xi\)
and hence:
\[
    \ext{\bm\beta}^\top\bm W\ext{\bm r}
    =(\bm Y\ext{\bm\beta})^\top\bm\Xi\ext{\bm r}
    =\tensorspan{\ext{\bm y}}^\top\bm\Xi\ext{\bm r}.
\]
Thus Eq.~\eqref{eq:pcs-query-matrix} checks the query-side
terminal value exactly. Completeness follows from
Fact~\ref{fact:ip} and the underlying PCS opening.
For \(m=1\), the initialization suffices, with no recurrence
or sumcheck rounds.

\smallskip\noindent\textbf{Soundness.}
Fix the matrix \(\bm A\) bound by the commitment and suppose that
Eq.~\eqref{eq:pcs-opening-relation} is false.
If the first scalar check accepts, then
\(\ext{\bm u}^{\top}\neq\ext{\bm\xi}^{\top}\bm A\), hence
\(\bm U\neq\bm\Xi^\top\bm A\).
Injectivity of basis packing makes
\((\bm U-\bm\Xi^\top\bm A)\ext{\bm\beta}\) a nonzero
column of \(\ExtensionField^\kappa\), fixed before
\(\ext{\bm r}\) is sampled. Its inner product with the uniform
\(\ext{\bm r}\) vanishes with probability
\(1/|\ExtensionField|\). Outside this event,
Eq.~\eqref{eq:pcs-random-opening} is false.

Except with probability \(2\ell/|\ExtensionField|\),
Fact~\ref{fact:ip} then leaves a false terminal claim in
Eq.~\eqref{eq:pcs-terminal-claims}.
The recurrence computes \(\bm W\) from the public tensor factors
and the terminal point, so a false query-side claim is rejected
deterministically. If that check passes, the witness-side claim
must be false. The underlying PCS accepts this adaptively chosen
false MLE claim with probability at most
\(\epsilon_{\mathsf{PCS}}\).
A union bound gives Eq.~\eqref{eq:pcs-soundness}.

\smallskip\noindent\textbf{Prover complexity.}
Expanding \(\ext{\bm\xi}\) and reading its basis coordinates
gives \(\bm\Xi\) in \(O(m)\) extension-field operations.
The matrix product \(\bm U=\bm\Xi^\top\bm A\) takes
\(\kappa N\) base-field multiplications. Its columns give the
coordinates of \(\ext{\bm u}^{\top}=\ext{\bm\beta}^{\top}\bm U\),
so packing requires only a change of representation.

Computing \(\bm\Xi\ext{\bm r}\) takes \(N\) mixed-field
multiplications; the extension-field inner-product reduction takes
\(4m+O(\ell)\) extension-field multiplications.
Packing the rows of \(\bm U\) and taking their inner product
with \(\ext{\bm r}\) computes \(\ext c\) using
\(\kappa\) extension-field multiplications.
The prover need not construct \(\bm Y\) or \(\bm W\).
These costs give Eq.~\eqref{eq:pcs-prover-cost}.
Under schoolbook arithmetic they are
\(O(\kappa N+\kappa^3)\,\bb\), which is
\(O(\kappa N)\,\bb\) when \(m\geq\kappa\).

The matrices \(\bm A,\bm\Xi\) occupy \(O(N)\) base-field
coordinates, and \(\bm U\) and the length-\(\kappa\)
extension-field vectors occupy \(O(\kappa^2)\).
A constant number of length-\(m\) extension-field buffers also
occupy \(O(m\kappa)=O(N)\) base-field coordinates.
Thus peak prover memory, excluding the underlying PCS workspace,
is \(O(N+\kappa^2)\), or \(O(N)\) when \(m\geq\kappa\).

\smallskip\noindent\textbf{Verifier complexity and communication.}
The verifier implements Eq.~\eqref{eq:pcs-query-recurrence}
without forming the multiplication matrices.
Right multiplication by \(\bm M_{\ext z}^\top\) scales each
row, interpreted in the fixed basis as an extension-field element,
by \(\ext z\). Left multiplication by \(\bm M_{\ext z}\)
scales each column in the same way. Each scaling therefore costs
\(\kappa\) extension-field multiplications.
The two terms in each recurrence step require three such scalings,
followed by \(\kappa^2\) base-field additions.
Over \(\ell\) steps, the cost is at most
\(3\kappa\ell\) extension-field multiplications and
\(O(\kappa^2\ell)\) base-field operations, including coordinate
rearrangement. A constant number of intermediate matrices use
\(O(\kappa^2)\) additional base-field coordinates.

The initial scalar check costs \(\kappa\) extension-field
multiplications. To compute \(\ext c\) and the expected terminal
value \(\ext w\), the verifier packs the rows of \(\bm U\) and
the columns of \(\bm W\), respectively, and takes one
length-\(\kappa\) extension-field inner product per target.
Packing is coordinate reinterpretation, so these targets
cost \(O(\kappa)\) extension-field multiplications and
\(O(\kappa^2)\) coordinate operations in total.
The inner-product verifier costs \(O(1+\ell)\) extension-field
operations, including its terminal scalar check.
This gives Eq.~\eqref{eq:pcs-verifier-cost}.

The prover sends \(\kappa\) extension-field elements for
\(\ext{\bm u}\), followed by the \(2\ell+2\) extension-field
elements of Fact~\ref{fact:ip}, including both terminal values.
The verifier computes \(\bm W\) locally, so no matrix message or
base-field proof is sent. Excluding the underlying opening, total
prover-to-verifier communication is therefore
\(\kappa+2\ell+2\) extension-field elements.
\end{proof}

\subsection{Proof of Lemma~\ref{eq:permutation}}
\label{appx:proof-copy-constraint}

\begin{proof}
The radix-\(\underline q\) representations of \(0,\ldots,L-1\)
use at most \(\lceil\log_{\underline q}L\rceil\) digits.
The digit injection preserves their distinctness.
Since \(\theta_1,\ldots,\theta_d\) are linearly independent
over \(\Field\), and the last coordinate is reserved for the
witness value:
\[
    a_i=b_j
    \quad\Longleftrightarrow\quad
    \bigl(\sigma(i)=j\ \text{and}\ w_i=w_j\bigr).
\]
Every label occurs exactly once in each operand.
Thus equality of the multisets forces
\(a_i=b_{\sigma(i)}\), and hence \(w_i=w_{\sigma(i)}\),
for all \(i\in[L]\).
Conversely, these copy equalities give
\(a_i=b_{\sigma(i)}\), proving the permutation relation.
The encoding is deterministic and injective, so it adds no
soundness error. Appending the same zeros to both operands
preserves the equivalence by multiset cancellation.
\end{proof}

\subsection{Proof of Proposition~\ref{proposition:position-encoding-evaluation}}
\label{appx:proof-position-encoding}

\begin{proof}
For the chosen digit injection, each label is the sum of the
contributions of its nonzero position bits.
The recursion and interval decomposition in
Appendix~\ref{appx:mleiota} therefore give the ordinary MLE
evaluation in \(O(\log L)\) operations.

For a packed query, split the position bits into column bits
and row bits in column-major order.
The unpadded label is the sum of a column contribution
\(\bm\iota_{\mathsf{col}}\) and a row contribution
\(\bm\iota_{\mathsf{row}}\).
For arbitrary column weights \(\ext{\bm r}\) and tensor row
weights \(\ext{\bm l}\), the full-table contraction is:
\[
    \left\langle\bm\iota_{\mathsf{col}},\ext{\bm r}\right\rangle
       \left(\bm1^\top\ext{\bm l}\right)
    +
    \left(\bm1^\top\ext{\bm r}\right)
       \left\langle\bm\iota_{\mathsf{row}},\ext{\bm l}\right\rangle.
\]
Recursive doubling generates all \(\kappa_{\mathbb K}\)
column contributions in \(O(\kappa_{\mathbb K})\) additions.
For each row-bit suffix, compute both its total tensor weight
and its label-weighted sum.
Splitting off one two-entry tensor factor updates both
quantities in constant time, including when its entries do
not sum to one.

Zero padding keeps the first \(L\) positions.
In column-major order, these consist of complete columns
and at most one partially filled column.
All complete columns use the same two full-row contractions.
For the partial column, the binary interval decomposition
computes the two corresponding prefix contractions using the
precomputed suffix quantities.
One scan of the prefix length costs \(O(\log L)\) operations.
Combining the column contributions then costs
\(O(\kappa_{\mathbb K})\) operations.
The total is \(O(\kappa_{\mathbb K}+\log L)\), with the same
bound on working field elements and no division.
\end{proof}

\subsection{Proof of Theorem~\ref{theorem:single-instance-packing}}
\label{appx:proof-single-instance}

\begin{proof}
\noindent\textbf{Completeness.}
For a satisfying assignment, the five scalar claims in
Eq.~\eqref{eq:plonkWeightedHadamard1} and
Eq.~\eqref{eq:plonkWeightedHadamard2} sum to zero.
Four packed Hadamard proofs reduce the selector--wire product
claims to evaluations of the selectors and
\(\bm w_L,\bm w_R,\bm w_M,\bm w_O\).
The fifth authenticates the remaining evaluation of
\(\bm w_M=\bm w_L\circ\bm w_R\).
All terminal queries have arbitrary column weights and tensor
row weights, as required by Section~\ref{sec:pcs}.

Lemma~\ref{eq:permutation} proves the copy equivalence.
The two operands have length \(L=3N+n_{\mathsf{io}}\leq4N\);
append zeros to reach length \(4N\).
Their grand products then share the padding factor
\(\ext z^{4N-L}\), which is nonzero because
\(\ext z\sample\ExtensionField\setminus\{0\}\).
The two grand-product reductions therefore prove exactly the
original permutation relation.
Their packing width \(\kappa_{\mathbb K}=\kappa/d\) and
row count \(4N/\kappa_{\mathbb K}=4md\) are powers of two,
with \(4md\geq\kappa_{\mathbb K}\geq2\).
Thus the parameters satisfy the interface of
Section~\ref{sec:permutation}, and every honest check accepts.

\smallskip\noindent\textbf{Terminal copy queries.}
Each copy terminal is a weighted evaluation of a padded operand,
with \(\kappa_{\mathbb K}\) arbitrary column weights and
\(4N/\kappa_{\mathbb K}\) tensor row weights.
By linearity, the two terminals involve
\(\bm\varsigma,\bm\iota\), and \(\bm w\).
Restricting to one length-\(N\) wire segment fixes the first
two position bits.
The last \(\log_2m\) bits remain tensor row factors.
Absorb all other free factors into the column weights to obtain
a length-\(\kappa\) column query on the original
\(m\times\kappa\) wire or digit matrix.
This changes the query representation, not the committed data.

Proposition~\ref{proposition:position-encoding-evaluation}
computes the identity-label contribution in
\(O(\kappa_{\mathbb K}+\log L)\) extension-field operations.
The verifier evaluates the public-I/O values and their permuted
labels directly. Traversing their contiguous interval and
reusing tensor prefixes costs
\(O(\kappa_{\mathbb K}+n_{\mathsf{io}}+\log L)\) operations.
No commitment to either contribution is needed.

\smallskip\noindent\textbf{Batched terminal openings.}
Write \(h=\lceil\log_{\underline q}L\rceil\).
The three fixed-basis wire packings form one witness vector
of logical length \(3m\), committed before any challenge.
The public vector concatenates the five selector packings
and the three wire-segment restrictions of each of the \(h\)
permuted-position digit tables.
Its logical length is \(m_{\mathsf{pub}}=(5+3h)m\).
The public commitment, circuit tables, and packing polynomials
are preprocessed.

The gate argument leaves five selector and five wire queries.
The two copy terminals add six wire queries and \(3h\)
permuted-digit queries, giving \(16+3h\) packed queries.
Apply the column-projection step of Section~\ref{sec:pcs}
to each query and check its claimed scalar value.
Fix all column messages before sampling independent projection
vectors for the different queries.
Within each group, zero-extend the projected queries to their
committed blocks and sum the resulting inner-product claims.
One invocation of Fact~\ref{fact:ip} per group reduces their
sum to one MLE opening of the concatenated vector.
The verifier evaluates the combined public query exactly,
using the recurrence of Section~\ref{sec:pcs} on each block
and retaining its block-prefix weight.
Thus there are two underlying PCS openings.

If a column message is false, its coordinate error satisfies
\((\bm U-\bm\Xi^\top\bm A)\ext{\bm\beta}\neq\bm0\).
Consequently, a group containing a false claim has a nonzero
linear error polynomial in its independent projection vectors.
It vanishes with probability at most
\(1/|\ExtensionField|\) per group.
The two inner-product reductions add
\(O(\log N)/|\ExtensionField|\) error.
Here the padded group lengths have logarithm \(O(\log N)\),
since \(h<d\leq\kappa/2\).

\smallskip\noindent\textbf{Soundness.}
Condition on binding of the witness commitment.
Its fixed-basis packing determines all three wire vectors
and the virtual product before any challenge.
The selectors and position tables are fixed at setup.

A nonzero gate residual vanishes at the random MLE point with
probability at most \(\log_2N/|\ExtensionField|\).
Otherwise an accepting scalar sum contains a false claim.
Acceptance then requires an error in one of the five Hadamard
reductions or in terminal authentication.
For the nested product claim, its value is fixed before the
fresh challenges of the fifth reduction.

If a copy constraint fails, the unpadded fingerprint difference
is nonzero and has degree at most \(L-1\).
At the nonzero challenge \(\ext z\), its padded version vanishes
if and only if the unpadded difference does.
The fingerprint error is therefore at most
\((L-1)/(|\ExtensionField|-1)\).
Let \(\epsilon_{\mathsf{GP}}\) be the bound of
Proposition~\ref{proposition:permutation-complexity},
instantiated with length \(4N\), packing width
\(\kappa_{\mathbb K}\), and base field \(\mathbb K\).
A union bound gives:
\[
\begin{aligned}
    \epsilon_{\mathsf{single}}\leq{}&
    \epsilon_{\mathsf{PCS}}
    +\frac{L-1}{|\ExtensionField|-1}
    +2\epsilon_{\mathsf{GP}}\\
    &+5\left(
        \frac{2\kappa-2}{|\ExtensionField|-\kappa}
        +\frac{2\log_2N}{|\ExtensionField|}
      \right)
    +O\left(\frac{\log N}{|\ExtensionField|}\right).
\end{aligned}
\]
The last term includes the initial gate projection and the
two batched opening reductions.
The grand-product contribution is
\(O((\kappa_{\mathbb K}\log_2\kappa_{\mathbb K}
+(\log_2N)^2)/|\ExtensionField|)\).
Since \(N=m\kappa\geq\kappa^2\), \(L\leq4N\), and
\((\log_2N)^2=O(N)\), this proves
Eq.~\eqref{eq:together-soundness}.
The term \(\epsilon_{\mathsf{PCS}}\) includes joint binding
and the two adaptive underlying openings.

\smallskip\noindent\textbf{Prover time.}
The five Hadamard reductions cost \(O(N\kappa)\bb\).
For \(q\geq2\kappa-1\), interpolation uses base-field nodes.
If \(\kappa<q<2\kappa-1\), the intermediate-field variant
uses degree-two interpolation-field tables; its online cost
is unchanged and its table storage is \(O(q^4)=O(\kappa^4)\).
The power-of-two \(\kappa\) permits this degree-two subfield.
The \(16+3h\) column projections and the two grouped
inner-product reductions cost \(O(N\kappa h)\bb\).

Choose polynomial bases along \(\Field\subset\mathbb K
\subset\ExtensionField\) for the copy computation.
A \(\mathbb K\)-multiplication costs \(O(d^2)\bb\).
The two grand-product reductions therefore cost:
\[
    O\left(
       N\kappa_{\mathbb K}\log_2\kappa_{\mathbb K}
       +\kappa_{\mathbb K}^4
    \right)d^2\bb
    =
    O\left(
       N\kappa d\log_2\kappa_{\mathbb K}
       +\frac{\kappa^4}{d^2}
    \right)\bb.
\]
The \(\log_2\kappa_{\mathbb K}\) factor counts the packed
coefficient levels; evaluating their aggregated polynomials
accounts for the additive term.
The least-power-of-two choice gives
\(h+1\leq d<2(h+1)\).
Since \(\log\underline q=\Theta(\log q)\) and \(L=\Theta(N)\),
both \(h\) and \(d\) are \(\Theta(\log_qN)\).
Combining these contributions gives the stated
\(O(N\kappa\log_qN\log_2\kappa)\bb+O(\kappa^4)\) bound.
The copy computation stays in its tower basis; only interface
queries and messages are converted to the fixed commitment basis.
These conversions cost \(O(\kappa^2)\bb\) per extension-field
element and are covered by the same bounds.

\smallskip\noindent\textbf{Verifier time and communication.}
The \(O(h)\) packed column checks and combined query evaluations
cost \(O(\kappa h\log_2N)\ee\).
The two grand-product reductions add
\(O(\kappa_{\mathbb K}\log_2\kappa_{\mathbb K}
+(\log_2N)^2)\ee\).
Direct identity-label and public-I/O evaluations contribute
\(O(\kappa_{\mathbb K}+n_{\mathsf{io}}+\log N)\ee\).
These costs give the verifier bound.

The packed column messages contain \(O(\kappa h)\)
extension-field elements.
The two grand-product transcripts contain
\(O(\kappa_{\mathbb K}\log_2\kappa_{\mathbb K}
+(\log_2N)^2)\) such elements.
The remaining Hadamard and batched inner-product transcripts
are absorbed by these terms.
Using \(h=O(\log_qN)\) and
\(\kappa_{\mathbb K}\leq\kappa\) gives the stated proof size.

\smallskip\noindent\textbf{Prover memory and PCS costs.}
The grand-product workspace uses
\(O(dN+d\kappa_{\mathbb K}^3)
=O(dN+\kappa^3/d^2)\) base-field coordinates.
The witnesses, circuit tables, and grouped opening buffers
occupy \(O(Nh)\) coordinates.
Sequential execution reuses the gate and copy workspaces.
Including the optional interpolation tables and public I/O
gives \(O(N\log_qN)\Field+O(\kappa^4+n_{\mathsf{io}})\).

The four PCS functions bound aggregate resources for one online
witness commitment and one opening of each of the two vectors,
of logical lengths \(3m\) and \(m_{\mathsf{pub}}\).
They include domain padding and both retained commitment states;
public commitment generation is preprocessed.
Their total argument is therefore \(3m+m_{\mathsf{pub}}\).
This proves Eq.~\eqref{eq:together-single-efficiency}.
Unlabeled time and memory remainders count base-field operations
and base-field elements, respectively.
\end{proof}

\subsection{Proof of Theorem~\ref{theorem:multiple-instance-packing}}
\label{appx:proof-multiple-instance}

\begin{proof}
\noindent\textbf{Completeness.}
Let \(\bm R\in\Field^{N\times\kappa}\) contain the gate
residuals of the \(\kappa\) instances as columns.
The left-hand side of Eq.~\eqref{eq:together-multiple-gates}
is \(\bm R\ext{\bm\beta}\), including the term
\(\bm q_C\bm1_\kappa^\top\ext{\bm\beta}\).
Linear independence of the basis gives
\(\bm R\ext{\bm\beta}=\bm0_N\) if and only if
\(\bm R=\bm0_{N\times\kappa}\).
Similarly, packing commutes with the common row permutation,
so \(\sigma(\ext{\bm w})=\ext{\bm w}\) is equivalent to
the copy constraints in every instance.
This equivalence includes the packed public I/O.
Eq.~\eqref{eq:together-multiple-product-opening} authenticates
the virtual product vector from
Eq.~\eqref{eq:together-multiple-wires}.
Thus all scalar checks and tree relations accept for an honest prover.

\smallskip\noindent\textbf{Public packing.}
Let \(h=\lceil\log_{\underline q}L\rceil\).
Form \(\bm Q\in\Field^{N\times\kappa}\) with first five columns
\(\bm q_L,\bm q_R,\bm q_M,\bm q_O,\bm q_C\).
The next three blocks of \(h\) columns contain the position digits
\(\bm\sigma_{(1)},\ldots,\bm\sigma_{(h)}\), restricted to the
three length-\(N\) wire segments, respectively.
Fill the remaining columns with zeros.
The capacity condition \(5+3h\leq\kappa\) therefore gives:
\[
    \bm Q\ext{\bm\beta}\in\ExtensionField^N.
\]
A selector query uses its coordinate unit vector as the column
weight in Eq.~\eqref{eq:pcs-opening-relation}.
For the permuted-label query on segment \(\ell\in[3]\), use
\(\ext\rho_{5+(\ell-1)h+j}=\ext\beta_j\) for \(j\in[h]\),
and zero elsewhere.
Thus five selector queries and three permuted-label queries
are weighted evaluations of the same public matrix.
The row weights are tensor vectors, so Section~\ref{sec:pcs}
applies to all eight queries.
The public commitment is preprocessed.
The identity-label MLE is evaluated directly by
Proposition~\ref{proposition:position-encoding-evaluation};
it is not committed.

\smallskip\noindent\textbf{Gate soundness.}
Condition on binding of the witness commitment.
Its three length-\(N\) blocks have unique matrices of
\(\Field\)-coordinates in the fixed basis.
Thus the gate residual and \(\ext{\bm w}_M\) are determined
by the committed witnesses.
If some gate is invalid, the packed residual is nonzero.
Its random length-\(N\) tensor projection vanishes with
probability at most \(\log_2N/|\ExtensionField|\).

For a fixed projection point \(\ext{\bm x}\), the gate
sumcheck polynomial is the product of the multilinear equality
weight with the residual expression formed from selector and wire
MLEs. Each selector--wire product has individual degree two,
and multiplication by the equality weight raises this to three.
The constant-selector term has lower degree, giving sumcheck
error at most \(3\log_2N/|\ExtensionField|\).
The remaining virtual-product evaluation is checked by one
packed Hadamard proof on matrices of size \(N\times\kappa\).
Since the total number of entries is \(N\kappa\), its
reduction error is bounded by:
\[
    \frac{2\kappa-2}{|\ExtensionField|-\kappa}
       +\frac{2\log_2(N\kappa)}{|\ExtensionField|}.
\]
All terminal queries are authenticated by the batched openings below.
The bound applies jointly to all \(\kappa\) instances because
packing is injective.

\smallskip\noindent\textbf{Copy fingerprint.}
Suppose that the packed assignment violates a copy constraint.
Consider the bivariate difference polynomial:
\[
    P(Y,Z)=
       \prod_{i=1}^{L}(Z+\iota_i+Y\ext w_i)
       -\prod_{i=1}^{L}(Z+\varsigma_i+Y\ext w_i).
\]
If \(P\) were zero, unique factorization in
\(\ExtensionField[Y,Z]\) would equate the two multisets of
linear factors, each monic in \(Z\).
Since the labels are distinct and
\(\varsigma_i=\iota_{\sigma(i)}\), matching the constant
coefficients pairs position \(i\) with \(\sigma^{-1}(i)\).
Matching the \(Y\)-coefficients then forces
\(\ext w_i=\ext w_{\sigma^{-1}(i)}\), contradicting the copy
violation. Hence \(P\ne0\).

Both products have homogeneous degree-\(L\) part
\(\prod_{i=1}^{L}(Z+Y\ext w_i)\), so \(P\) has total degree
at most \(L-1\).
The independent uniform challenges \(\ext y,\ext z\)
therefore satisfy \(P(\ext y,\ext z)=0\) with probability at
most \(\min\{1,(L-1)/|\ExtensionField|\}\).
Padding each product-tree input with ones does not change
this fingerprint.

\smallskip\noindent\textbf{Product-tree soundness.}
Set \(\mu=\lceil\log_2L\rceil=\log_2(4N)\).
For one grand product, unfold the binary tree from its root
to the padded length-\(4N\) input vector.
At a level with \(2^\ell\) parent entries, where
\(\ell\in\{0,\ldots,\mu-1\}\), let
\(\ext{\bm x}\in\ExtensionField^\ell\) be the current
evaluation point.
The weighted product identity uses the polynomial:
\[
    \left\langle
       \tensorspan{\ext{\bm x}},\tensorspan{\bm X}
    \right\rangle
    \left\langle\ext{\bm c}_U,\tensorspan{\bm X}\right\rangle
    \left\langle\ext{\bm c}_D,\tensorspan{\bm X}\right\rangle,
\]
where \(\ext{\bm c}_U,\ext{\bm c}_D\) are the two child subvectors
and \(\bm X\) denotes \(\ell\) formal variables.
Its degree in each variable is at most three.
Sumcheck reduces its Boolean sum to child evaluations
\(\ext v_U,\ext v_D\) at one point \(\ext{\bm y}\), with
error at most \(3\ell/|\ExtensionField|\).
After both values are fixed, sample a fresh
\(\ext y_0\sample\ExtensionField\) and form the next claim:
\[
    (1-\ext y_0)\ext v_U+\ext y_0\ext v_D
       =
    \left\langle
       \ext{\bm c}_U\mathbin{\|}\ext{\bm c}_D,
       \tensorspan{(\ext y_0,\ext{\bm y})}
    \right\rangle.
\]
If either child claim is false, this merge accepts it with
probability at most \(1/|\ExtensionField|\).
Thus the two trees have total reduction error at most:
\[
    2\sum_{\ell=0}^{\mu-1}\frac{3\ell+1}{|\ExtensionField|}
       =\frac{3\mu(\mu-1)+2\mu}{|\ExtensionField|}.
\]
The \(\ell=0\) level has no sumcheck rounds but still uses
a fresh merge challenge.

Each tree ends in a single MLE claim on its padded input.
At its terminal point, the known shift and padding contribute:
\[
    \ext z\sum_{i=1}^{L}
       \bigl(\tensorspan{\ext{\bm x}}\bigr)_i
    +\sum_{i=L+1}^{4N}
       \bigl(\tensorspan{\ext{\bm x}}\bigr)_i
    =
    1+(\ext z-1)\sum_{i=1}^{L}
       \bigl(\tensorspan{\ext{\bm x}}\bigr)_i.
\]
The prefix sum takes \(O(\log N)\) operations by a binary
interval decomposition.
Each length-\(N\) wire segment fixes two prefix bits.
Retain their tensor weight and query the suffix MLE.
The public-I/O values and permuted labels at the remaining
\(n_{\mathsf{io}}\) positions are evaluated directly.
The identity-label MLE is computed by Appendix~\ref{appx:mleiota}.
Consequently, the two copy terminals leave six ordinary wire
queries and three weighted queries to the public matrix.

\smallskip\noindent\textbf{Batched openings and composition.}
There are nine ordinary wire queries: three from the gate check
and six from the copy check.
The two Hadamard terminals and eight public queries give
ten packed queries.
Apply the column-projection step of Section~\ref{sec:pcs}
to each packed query.
Fix all column messages and ordinary claimed values before
sampling independent projection vectors for the packed queries
and independent scalar weights for the ordinary queries.
Extend each query by zeros outside its committed block, and
sum the resulting inner-product claims within each group.

If an ordinary claim is false, its scalar weight has a nonzero
error coefficient.
If a packed column message is false and its scalar check passes,
injectivity gives
\((\bm U-\bm\Xi^\top\bm A)\ext{\bm\beta}\ne\bm0\).
Thus a group containing a false claim has a nonzero linear
error polynomial in its independent batching challenges.
Cancellation occurs with probability at most
\(1/|\ExtensionField|\) per group.

The witness concatenation has logical length \(3N\), padded
to \(4N\); the public packing has length \(N\).
One invocation of Fact~\ref{fact:ip} per group therefore adds
\(2\log_2(4N)/|\ExtensionField|\) and
\(2\log_2N/|\ExtensionField|\) error, respectively.
The verifier computes the combined query-side evaluation in
each group exactly: use the recurrence of Section~\ref{sec:pcs}
for packed queries and tensor contractions for ordinary
queries, retaining the witness block-prefix weights.
This leaves two underlying PCS openings.
Their joint binding and opening error is
\(\epsilon_{\mathsf{PCS}}\).

Combining these bounds gives:
\[
\begin{aligned}
    \epsilon_{\mathsf{multi}}
    \leq{}&\epsilon_{\mathsf{PCS}}
       +\frac{2\kappa-2}{|\ExtensionField|-\kappa}\\
       &+\frac{
          L-1+3\mu(\mu-1)+2\mu
          +8\log_2N+2\log_2(N\kappa)+6
       }{|\ExtensionField|}.
\end{aligned}
\]
Here the constant \(6\) includes the two aggregation checks
and the extra prefix rounds for the length-\(4N\) witness.
Since \(L\leq4N\), \(\kappa\leq N\), and
\((\log_2N)^2=O(N)\), this proves
Eq.~\eqref{eq:together-multiple-soundness}.

\smallskip\noindent\textbf{Efficiency.}
Each wire matrix has \(N\kappa\) entries.
Computing \(\bm W_L\circ\bm W_R\) costs \(O(N\kappa)\bb\).
The Hadamard reduction and ten packed column projections each
cost \(O(N\kappa^2+\kappa^3)\bb\), which is
\(O(N\kappa^2)\bb\) because \(N\geq\kappa\).
The assumption \(q\geq2\kappa-1\) supplies the interpolation
nodes over \(\Field\), so no intermediate-field table is needed.
The gate sumcheck and the two final inner-product reductions
use \(O(N)\) extension-field operations.
The product-tree layer lengths sum to \(O(N)\), and their
sumchecks cost \(O(N+\mu^2)\) extension-field operations.
Under schoolbook arithmetic, these contributions give
\(O(N\kappa^2)\bb\) total algebraic prover time.

The packed Hadamard check uses \(O(\kappa)\ee\) verifier time,
and the ten column projections and terminal query contractions
use \(O(\kappa\log_2N)\ee\).
The gate, tree, and final inner-product checks add
\(O((\log_2N)^2)\ee\).
Packing the public I/O reads \(O(\kappa n_{\mathsf{io}})\)
base-field coordinates, and its terminal contractions and
the direct identity-label evaluations take
\(O(n_{\mathsf{io}}+\log N)\) extension-field operations.
This gives the stated verifier time.

The ten column messages contain \(O(\kappa)\)
extension-field elements in total.
The Hadamard proof, gate sumcheck, tree reductions, and final
inner-product proofs add
\(O(\kappa+(\log_2N)^2)\) extension-field elements,
giving the proof-size bound.
The witnesses and public matrix occupy \(O(N\kappa)\)
base-field coordinates.
Sequential execution retains \(O(N)\) extension-field
tree and sumcheck entries and \(O(\kappa^2)\) coordinates
for the column messages.
Thus algebraic prover memory is \(O(N\kappa)\Field\).

The online PCS costs cover one witness commitment and two
openings on vectors of logical lengths \(3N\) and \(N\).
Their aggregate argument is therefore \(4N\), with the required
witness padding included in the PCS cost functions.
Public commitment generation is part of circuit preprocessing;
both retained commitment states count toward PCS memory.
These calculations prove
Eq.~\eqref{eq:together-multiple-efficiency}.
Dividing the algebraic prover time and peak memory by
\(\kappa\) gives the amortized per-instance bounds.
\end{proof}
